\section{Functional Evidence}
\label{sec:functional}

\lstdefinestyle{run}{basicstyle=\ttfamily\scriptsize,columns=fixed,breaklines=true,frame=single,backgroundcolor=\color{codebg}}

Everything in this section is copied from transcripts produced by the submitted
Stage~3 build. We did not edit the console output, apart from cutting long stretches,
which are marked with \texttt{[...]}. Each run is stored in the repository under
\texttt{test-runs/} as a pair of files. The \texttt{.txt} file holds the exact keyboard
input: the stage, the seed, then one line per choice. The \texttt{.out} file is the
transcript we recorded. A TA can replay any run with:

\begin{lstlisting}[style=run]
javac -d bin src/restaurantrush/*.java
java -cp bin restaurantrush.RestaurantRush --echo < test-runs/stage1-victory.txt
\end{lstlisting}

\texttt{--echo} prints each input after its prompt, so the replay reads like a live
game. If an input file ends early, every remaining prompt takes its last option (Wait,
Back or ``n''), so the game still plays all 18 turns. \texttt{TESTING.md} in the
repository describes each run in more detail.

\subsection{Summary of reproducible runs}

All runs use the base setting from the brief: 18 turns, 2 tables, one Waiter, one Chef
and one Cashier, with arrivals on turns 1, 4, 7, 10, 13 and 16. Stage~3 adds the Host.
The ``victory'' runs follow one simple rule each turn. The Waiter serves a READY dish if
there is one; otherwise takes an order; otherwise seats someone. The Chef always works on
the oldest unfinished order. The ``defeat'' runs use the same Waiter, but the Chef never
cooks.

\begin{table}[H]
\centering
\small
\caption{Reproducible runs (all numbers taken from the final screen of each \texttt{.out} file)}
\label{tab:runs}
\begin{tabular}{@{}c c l l c r c@{}}
\toprule
Stage & Seed & Input file (\texttt{test-runs/}) & Result & Paid & Revenue & Avg sat. \\
\midrule
1 & 42 & \texttt{stage1-victory.txt}                  & Victory & 5 & \$63.00 & 74.4 \\
1 & 42 & \texttt{stage1-defeat.txt}                   & Defeat  & 0 & \$0.00  & N/A  \\
1 & 42 & \texttt{stage1-invalid-input.txt}            & Defeat  & 1 & \$18.00 & 76.0 \\
2 & 42 & \texttt{stage2-victory.txt}                  & Victory & 4 & \$53.10 & 76.3 \\
2 & 42 & \texttt{stage2-defeat.txt}                   & Defeat  & 0 & \$0.00  & N/A  \\
2 & 42 & \texttt{extra-critic-slow-service.txt}       & Victory & 4 & \$53.10 & 76.3 \\
3 & 42 & \texttt{stage3-victory.txt}                  & Victory & 5 & \$88.45 & 83.2 \\
3 & 42 & \texttt{stage3-defeat.txt}                   & Defeat  & 0 & \$0.00  & N/A  \\
3 & 42 & \texttt{stage3-happy-hour-after-payment.txt} & Victory & 5 & \$88.74 & 87.6 \\
3 & 1  & \texttt{stage3-vip-happy-hour-seed1.txt}     & Victory & 5 & \$73.80 & 91.4 \\
3 & 42 & \texttt{extra-host-and-waiter-seat.txt}      & Defeat  & 0 & \$0.00  & N/A  \\
\bottomrule
\end{tabular}
\end{table}

The \texttt{invalid-input} and \texttt{extra-*} runs exist to show one specific behaviour.
They stop giving real choices after a few turns, so their final results are not
meaningful. No victory run has more than five paid customers. This matches the note in
the brief: the customer who arrives on turn 16 cannot pay within 18 turns.

\subsection{Start-of-turn display}

The brief asks for the turn number, revenue, the average satisfaction of paid customers,
the queue, the occupied tables, each customer's satisfaction, order status, and the staff
tasks. All of these appear at the top of every turn. In Stage~3, the Happy Hour status
and the Host are also shown. Any Happy Hour recovery is applied first, so the table
shows the satisfaction after the +15 (\texttt{stage3-victory.out}, turn 8):

\begin{lstlisting}[style=run]
========================================================================
 TURN 8 / 18     Revenue: $9.00     Avg satisfaction (paid): 92.0
========================================================================
  * Happy Hour: VIP #3 cheers up, satisfaction 95 -> 100.
 Happy Hour: ACTIVE, 1 turn(s) left including this one (20% off new orders, waiting customers +15)
 Waiting queue:
   (empty)
 Tables:
   Table 1: VIP #3 [satisfaction 100, ORDERED]
            order: COOKING 1/2 units, $15.55
   Table 2: Regular #2 [satisfaction 84, READY_TO_PAY]
            order: READY_TO_PAY 2/2 units, $21.60
 Staff on duty:
   Host Hana        seat one waiting customer (in addition to the Waiter)
   Waiter Wendy     one task: seat, take an order or serve a READY dish
   Chef Carlos      one preparation unit on one order
   Cashier Cathy    automatic, collects one payment per turn
\end{lstlisting}

At the end of each turn, the game prints the served count, unhappy departures, revenue,
average satisfaction and turns used:

\begin{lstlisting}[style=run]
 End of turn 8: served 2 | unhappy departures 0 | revenue $30.60
                 avg satisfaction (paid) 88.0 | turns used 8/18
\end{lstlisting}

\subsection{Invalid choices are explained and do not use up the task}

In \texttt{stage1-invalid-input.out}, the manager types text, then an out-of-range
number, then two actions that are impossible on turn 1. The Waiter's menu reappears
every time, and the turn continues only after a valid action:

\begin{lstlisting}[style=run]
 Waiter Wendy - choose one task:
   1. Seat a waiting customer
   2. Take an order
   3. Serve a READY dish
   4. Wait
 Choice > abc
  ! Please enter a number from 1 to 4.
 Choice > 9
  ! Please enter a number from 1 to 4.
 Choice > 3
  ! Not possible: There are no orders in the kitchen. Choose again.
 Waiter Wendy - choose one task:
   [...]
 Choice > 2
  ! Not possible: No customer is seated yet. Choose again.
\end{lstlisting}

The reasons come from the \texttt{ActionResult} returned by \texttt{Waiter}, not from the
view. The same run also tries to serve a half-cooked Pasta and to seat a customer at an
occupied table:

\begin{lstlisting}[style=run]
 Serve which dish?
   1. Pasta for Regular #1 - COOKING 1/2 units, $18.00
   2. Back
 Choice > 1
  ! Not possible: Pasta for Regular #1 is not READY yet (1/2 units done). Choose again.
 [...]
  ! Not possible: Table 1 is occupied by Regular #1. Choose again.
\end{lstlisting}

\subsection{READY, SERVED, READY\_TO\_PAY, PAID: the worked example from the brief}

The first six turns of \texttt{stage1-victory.out} repeat the example in Section~3.3 of
the brief. Regular \#1 orders Pasta (2 units) and gets the manager's full attention. The
customer's satisfaction at the end of turns 1--3 is 92, 84 and 76. It then stays at 76
once the dish is served, and 76 is the value recorded at payment.

\begin{lstlisting}[style=run]
 TURN 2  * Waiter Wendy takes Regular #1's order: Pasta ($18.00).
         * Chef Carlos works on Pasta for Regular #1 (1/2).
 TURN 3  * Chef Carlos works on Pasta for Regular #1 (2/2) - READY, can be served from next turn.
 TURN 4    Table 1: Regular #1 [satisfaction 76, ORDERED]
                    order: READY 2/2 units, $18.00
         * Waiter Wendy serves Pasta for Regular #1.
 TURN 5    Table 1: Regular #1 [satisfaction 76, EATING]
         * Regular #1 finishes eating and is READY_TO_PAY.
 TURN 6  * Cashier Cathy collects $18.00 from Regular #1 (satisfaction recorded: 76). Table freed.
\end{lstlisting}

To save space, the excerpt above shows only the relevant lines of turns 2--6; each
line is copied word for word. The ordering rules can all be seen here. The Chef cooks an
order the Waiter took in the same turn (turn 2). A dish that becomes READY on turn 3
cannot be served until turn 4. A customer who becomes READY\_TO\_PAY on turn 5 pays on
turn 6, at the earliest.

\subsection{Stage 2: VIP and Critic behaviour}

\textbf{VIP pricing.} VIP \#3 orders a Salad, and the 10\% discount is locked into the
order when it is taken (\texttt{stage2-victory.out}):

\begin{lstlisting}[style=run]
  * Waiter Wendy takes VIP #3's order: Salad (list $9.00, charged $8.10).
  [...]
  * Cashier Cathy collects $8.10 from VIP #3 (satisfaction recorded: 85). Table freed.
\end{lstlisting}

The VIP's satisfaction also drops more slowly: VIP \#3 lost 5 points per turn and was
recorded at 85, while Regular \#1 and Regular \#2 in the same run lost 8 per turn and were
recorded at 76.

\textbf{Critic preference and reaction.} We documented that a Critic always orders the
most expensive single dish (Pasta). A Critic gains 5 satisfaction if served within 2 turns
of ordering and loses 10 if it takes 3 or more turns. Both cases appear in the
transcripts:

\begin{lstlisting}[style=run]
stage2-victory.out (ordered turn 15, served turn 17):
  * Waiter Wendy takes Critic #5's order: Pasta ($18.00).
  [...]
  * Waiter Wendy serves Pasta for Critic #5. Critic #5 reacts: satisfaction 52 -> 57.

extra-critic-slow-service.out (same game, Waiter waits on turn 17, serves on turn 18):
  * Waiter Wendy serves Pasta for Critic #5. Critic #5 reacts: satisfaction 40 -> 30.
\end{lstlisting}

The Critic arrives on turn 13 and is served too late to pay, so the slow-service run
ends with the same result as the normal Stage~2 victory run. The Critic also
loses satisfaction faster (12 per turn): it arrived at 100 and was at 52 four turns later,
when the Pasta reached the table.

\subsection{Stage 3: Host, ComboMeal and Happy Hour}

\textbf{Host in addition to the Waiter.} In \texttt{extra-host-and-waiter-seat.out}, two
customers are queuing on turn 4. The Host seats one, and the Waiter seats the other in the
same turn:

\begin{lstlisting}[style=run]
 Host Hana - seat which customer?
   1. Regular #1 [satisfaction 76, WAITING]
   2. Regular #2 [satisfaction 100, WAITING]
   3. Wait
 Choice > 1
 [...]
  * Host Hana seats Regular #1 at Table 1.
 [...]
  * Waiter Wendy seats Regular #2 at Table 2.
\end{lstlisting}

\textbf{ComboMeal.} The Stage~3 menu lists all six combos. Each combo's price is 90\% of
its two parts, rounded to the nearest cent, and its preparation units are the sum of its
two parts. For example, $(\$9 + \$18) \times 0.9 = \$24.30$, with $1 + 2 = 3$ units.
Customers in Stage~3 order a random combo 30\% of the time, using the seeded generator.

\begin{lstlisting}[style=run]
  Salad + Salad Combo          $16.20   2 prep unit(s)
  Salad + Burger Combo         $21.60   2 prep unit(s)
  Salad + Pasta Combo          $24.30   3 prep unit(s)
  Burger + Burger Combo        $27.00   2 prep unit(s)
  Burger + Pasta Combo         $29.70   3 prep unit(s)
  Pasta + Pasta Combo          $32.40   4 prep unit(s)
  [...]
  * Waiter Wendy takes Regular #4's order: Salad + Pasta Combo ($24.30).
\end{lstlisting}

\textbf{Happy Hour.} The manager switches it on at the start of turn 7, before any staff
act. It lasts for that turn and the next. The remaining duration is printed at the top of
both turns: ``2 turn(s) left including this one'' on turn 7 and ``1 turn(s) left'' on
turn 8 (see the start-of-turn display above). With seed 1, the VIP orders a Burger during Happy Hour. The
VIP discount is applied first and Happy Hour second, giving exactly the brief's example of
$\$15 \times 0.9 \times 0.8 = \$10.80$:

\begin{lstlisting}[style=run]
stage3-vip-happy-hour-seed1.out:
  * Waiter Wendy takes VIP #3's order: Burger (list $15.00, charged $10.80).

stage3-victory.out (combo for a VIP: $21.60 x 0.9 = $19.44, x 0.8 = $15.552 -> $15.55):
  * Happy Hour is ON for this turn and the next: 20% off new orders.
  [...]
  * Waiter Wendy takes VIP #3's order: Salad + Burger Combo (list $21.60, charged $15.55).
\end{lstlisting}

\texttt{stage3-happy-hour-after-payment.out} checks the recovery rule. Happy Hour is
switched on at turn 13. By then three customers have already paid, with recorded
satisfaction 92, 84 and 90 (average 88.7). Only customers still waiting for a table or
for food recover 15 points, on the activation turn and on the following turn. The paid
average stays at 88.7 throughout (only the relevant lines of turns 13--15 are shown):

\begin{lstlisting}[style=run]
 TURN 13 / 18     Revenue: $50.04     Avg satisfaction (paid): 88.7
  * Happy Hour is ON for this turn and the next: 20% off new orders.
  * Happy Hour: Regular #4 cheers up, satisfaction 76 -> 91.
 Happy Hour: ACTIVE, 2 turn(s) left including this one (20% off new orders, waiting customers +15)
 TURN 14 / 18     Revenue: $50.04     Avg satisfaction (paid): 88.7
  * Happy Hour: Critic #5 cheers up, satisfaction 88 -> 100.
 Happy Hour: ACTIVE, 1 turn(s) left including this one (20% off new orders, waiting customers +15)
 TURN 15 / 18     Revenue: $50.04     Avg satisfaction (paid): 88.7
 Happy Hour: used
\end{lstlisting}

The final average of 87.6 equals $(92+84+90+91+81)/5$. These are the values recorded at
each payment, so Happy Hour never changes a customer who has already paid. Orders placed
before turn 13 keep their original prices. For example, Regular \#4's combo was ordered
at \$24.30 and was still charged \$24.30.

\subsection{Abandonment, Victory and Defeat}

In the defeat runs the Chef never cooks. Satisfaction keeps falling by 8 per turn, and
Regular \#1 leaves on turn 13 without paying. The leaving customer's order is cancelled
and their table is freed (\texttt{stage1-defeat.out}):

\begin{lstlisting}[style=run]
  * Chef Carlos waits this turn.
  * Regular #1 runs out of patience and leaves without paying.
------------------------------------------------------------------------
 End of turn 13: served 0 | unhappy departures 1 | revenue $0.00
                 avg satisfaction (paid) N/A | turns used 13/18
\end{lstlisting}

The game always plays all 18 turns and then checks the three targets:

\begin{lstlisting}[style=run]
stage1-victory.out:
  VICTORY! The restaurant hit every target.
  Paid customers    5        (need 4)   [met]
  Revenue           $63.00   (need $45.00)  [met]
  Avg satisfaction  74.4     (need 60)  [met]
  Unhappy departures: 0

stage1-defeat.out:
  DEFEAT. At least one target was missed.
  Paid customers    0        (need 4)   [missed]
  Revenue           $0.00    (need $45.00)  [missed]
  Avg satisfaction  N/A      (need 60)  [missed]
  Unhappy departures: 2
\end{lstlisting}

\textbf{The turn-16 customer.} The game says at the start that the customer arriving on
turn 16 cannot finish paying in time. In \texttt{stage1-victory.out}, the best case for
that customer is an order taken on turn 18 and a dish that becomes READY on the same turn.
The game then ends with 5 of the 6 customers paid, as the brief expects:

\begin{lstlisting}[style=run]
(The customer arriving on turn 16 cannot finish paying in time - that is expected.)
[...]
  * Waiter Wendy takes Regular #6's order: Burger ($15.00).
  [...]
  * Chef Carlos works on Burger for Regular #6 (1/1) - READY, can be served from next turn.
\end{lstlisting}
